\documentclass[12pt,letterpaper]{report}
\newcommand{\ProyectoCabecera}{PROYECTO: LÍNEA PRIMARIA LIMBANI - HUANCASAYANI, PROVINCIA DE SANDIA - DEPARTAMENTO DE PUNO}
\newcommand{\CursoCabecera}{\shortstack{INSTALACIONES\\ELÉCTRICAS II}}
\newcommand{\DocumentoPie}{ESPECIFICACIONES Y MONTAJE - LÍNEA PRIMARIA}
\newcommand{\AutorDocumento}{Aquiles Taylor Ramos Yapo}
\newcommand{\CursoCaratula}{INSTALACIONES ELÉCTRICAS II}
\newcommand{\TituloCaratula}{PROYECTO: “LÍNEA PRIMARIA LIMBANI - HUANCASAYANI”}
\newcommand{\UbicacionCaratula}{DISTRITO DE LIMBANI - PROVINCIA DE SANDIA - DEPARTAMENTO DE PUNO}
\newcommand{\TipoDocumentoCaratula}{ESPECIFICACIONES Y MONTAJE}
\newcommand{\FichaCaratula}{Estudiante: & Aquiles Taylor Ramos Yapo \\
Código: & 227368 \\
Docente del curso: & Por confirmar \\
Ubicación: & Limbani - Huancasayani, Sandia, Puno \\
Revisión técnica: & LP-R02 \ensuremath{\cdot} Presentación P01 \\
Edición de presentación: & 18 de octubre de 2026 \\
Naturaleza: & Proyecto académico estudiantil \\}
% Adaptación de proyectos/aquiles/expediente/preambulo.tex.
% Las variables de identidad y alcance se definen por documento.
\usepackage[utf8]{inputenc}
\usepackage[T1]{fontenc}
\usepackage[spanish,es-tabla]{babel}
\usepackage{geometry,amsmath,amssymb,graphicx,float,booktabs,tabularx,longtable,array,enumitem,xcolor}
\usepackage[scaled]{helvet}
\renewcommand{\familydefault}{\sfdefault}
\usepackage{hyperref,fancyhdr,titlesec,tocloft,tikz,colortbl,needspace,etoolbox}
\usepackage[normalem]{ulem}
\definecolor{modelblue}{RGB}{91,155,213}
\definecolor{headingred}{RGB}{220,0,0}
\definecolor{azul}{RGB}{91,155,213}
\definecolor{rojo}{RGB}{220,0,0}
\definecolor{gris}{RGB}{70,82,92}
\geometry{letterpaper,left=1.08in,right=.92in,top=1.55in,bottom=.92in,headheight=78pt,headsep=18pt,footskip=36pt}
\providecommand{\ProyectoCabecera}{PROYECTO ACADÉMICO}
\providecommand{\CursoCabecera}{INGENIERÍA}
\providecommand{\DocumentoPie}{MEMORIA ACADÉMICA}
\providecommand{\AutorDocumento}{Estudiante}
\providecommand{\LogoDocumento}{figuras/escudo-unap.png}
\hypersetup{colorlinks=true,linkcolor=black,citecolor=black,urlcolor=blue,pdfauthor={\AutorDocumento},pdftitle={\ProyectoCabecera}}
\urlstyle{same}
\setlength{\parindent}{0pt}
\setlength{\parskip}{.62em}
\setlength{\tabcolsep}{4pt}
\renewcommand{\arraystretch}{1.18}
\setlength{\headwidth}{\textwidth}
\setlist[itemize]{leftmargin=1.05cm,itemsep=.05em,topsep=.2em,label=\(\triangleright\)}
\setlist[enumerate]{leftmargin=1.05cm,itemsep=.05em,topsep=.2em}
\setcounter{secnumdepth}{3}\setcounter{tocdepth}{2}
\addto\captionsspanish{\renewcommand{\contentsname}{CONTENIDO}}
\renewcommand{\cfttoctitlefont}{\color{modelblue}\Large}
\renewcommand{\cftaftertoctitle}{\par\vspace{.4em}}
\setlength{\cftbeforetoctitleskip}{0pt}\setlength{\cftaftertoctitleskip}{.75em}
\renewcommand{\cftchapfont}{\bfseries}\renewcommand{\cftchappagefont}{\bfseries}
\renewcommand{\cftchapleader}{\cftdotfill{\cftdotsep}}
\renewcommand{\cftsecleader}{\cftdotfill{\cftdotsep}}
\renewcommand{\cftsubsecleader}{\cftdotfill{\cftdotsep}}
\setlength{\cftsecindent}{.25in}\setlength{\cftsubsecindent}{.60in}
\setlength{\cftbeforechapskip}{.30em}\setlength{\cftbeforesecskip}{.16em}
\titleformat{\chapter}[display]{\normalfont\bfseries\color{black}}{\normalsize\color{headingred}CAPÍTULO \thechapter}{6pt}{\LARGE}
\titleformat{name=\chapter,numberless}{\normalfont\bfseries\LARGE\color{black}}{}{0pt}{}
\titleformat{\section}{\normalfont\bfseries\Large\color{headingred}}{\thesection}{1em}{\uline}
\titleformat{\subsection}{\normalfont\bfseries\normalsize\color{headingred}}{\thesubsection}{1em}{}
\titleformat{\subsubsection}{\normalfont\bfseries\normalsize\color{headingred}}{\thesubsubsection}{1em}{}
\titlespacing*{\chapter}{0pt}{0pt}{.95em}
\titlespacing*{\section}{0pt}{1em}{.55em}
\titlespacing*{\subsection}{0pt}{.85em}{.35em}
\pretocmd{\section}{\Needspace{5\baselineskip}}{}{}
\pretocmd{\subsection}{\Needspace{4\baselineskip}}{}{}
\newcolumntype{L}[1]{>{\raggedright\arraybackslash}p{#1}}
\newcommand{\memoriaheader}{%
 \begingroup\arrayrulecolor{modelblue}\setlength{\arrayrulewidth}{1.1pt}\setlength{\tabcolsep}{4pt}\renewcommand{\arraystretch}{1.18}%
 \makebox[.90\headwidth][c]{\begin{minipage}{.90\headwidth}\centering
 \resizebox{\linewidth}{!}{\begin{tabular}{|>{\centering\arraybackslash}m{3.55in}|>{\centering\arraybackslash}m{1.05in}|>{\centering\arraybackslash}m{.55in}|}\hline
 \scriptsize\ProyectoCabecera & \scriptsize\CursoCabecera & \includegraphics[height=.50in]{\LogoDocumento}\\\hline
 \end{tabular}}\par\vspace{.13in}\textcolor{modelblue}{\rule{\linewidth}{.45pt}}
 \end{minipage}}\endgroup}
\newcommand{\memoriafooter}{\makebox[.90\headwidth][c]{\begin{minipage}{.90\headwidth}
 \textcolor{modelblue}{\rule{\linewidth}{.45pt}}\par\vspace{.08in}{\small\DocumentoPie\hfill\thepage}\end{minipage}}}
\fancypagestyle{memoria}{\fancyhf{}\fancyhead[C]{\makebox[0pt][c]{\memoriaheader}}\fancyfoot[C]{\makebox[0pt][c]{\memoriafooter}}\renewcommand{\headrulewidth}{0pt}\renewcommand{\footrulewidth}{0pt}}
\fancypagestyle{plain}{\fancyhf{}\fancyhead[C]{\makebox[0pt][c]{\memoriaheader}}\fancyfoot[C]{\makebox[0pt][c]{\memoriafooter}}\renewcommand{\headrulewidth}{0pt}\renewcommand{\footrulewidth}{0pt}}
\pagestyle{memoria}
\newcommand{\estado}[1]{\par\noindent\colorbox{modelblue!8}{\parbox{\dimexpr\linewidth-2\fboxsep}{\small\textbf{Condiciones del modelo.} #1}}\par}
\newcommand{\dato}[1]{\csname dato#1\endcsname}
\newcommand{\fuente}[1]{{\small\color{gris}\textit{Fuente / trazabilidad: #1}}\par}
\emergencystretch=3em
\makeatletter\setlength{\@fptop}{0pt}\makeatother
\sloppy

\renewcommand{\thesection}{\arabic{section}.}\renewcommand{\thesubsection}{\arabic{section}.\arabic{subsection}.}
\begin{document}\hypersetup{pageanchor=false}\pagenumbering{gobble}% Geometría y ornamentación de la carátula del expediente de vivienda de Aquiles.
\begin{titlepage}\thispagestyle{empty}
\begin{tikzpicture}[remember picture,overlay]
 \fill[modelblue!80!black] ([xshift=.92in,yshift=-.28in]current page.north west) rectangle ([xshift=1.08in,yshift=.28in]current page.south west);
 \draw[modelblue!80!black,line width=.7pt] ([xshift=1.10in,yshift=.28in]current page.south west) .. controls +(.10in,2.15in) and +(-.85in,-1.65in) .. ([xshift=2.55in,yshift=3.15in]current page.south west);
 \draw[green!45!black,line width=1pt] ([xshift=1.05in,yshift=.28in]current page.south west) .. controls +(.35in,2in) and +(-.95in,-1.55in) .. ([xshift=2.20in,yshift=3.85in]current page.south west);
 \draw[green!55!black,line width=1pt] ([xshift=1in,yshift=.28in]current page.south west) .. controls +(.20in,1.75in) and +(-.85in,-1.45in) .. ([xshift=1.70in,yshift=4.25in]current page.south west);
 \draw[modelblue!75,line width=.55pt] ([xshift=1in,yshift=.28in]current page.south west) .. controls +(.15in,1.50in) and +(-.80in,-1.35in) .. ([xshift=2.45in,yshift=2in]current page.south west);
\end{tikzpicture}
\vspace*{-.4in}\noindent\hspace*{1.20in}\begin{minipage}{5.20in}\centering
{\fontsize{13}{16}\selectfont\bfseries UNIVERSIDAD NACIONAL DEL ALTIPLANO\par}
{\fontsize{10}{12}\selectfont\bfseries FACULTAD DE INGENIERÍA MECÁNICA ELÉCTRICA, ELECTRÓNICA Y SISTEMAS\par}
{\fontsize{10}{12}\selectfont\bfseries ESCUELA PROFESIONAL DE INGENIERÍA MECÁNICA ELÉCTRICA\par}
\vspace{.20in}\includegraphics[height=1in]{\LogoDocumento}\par
\vspace{.20in}{\fontsize{12}{15}\selectfont\bfseries CURSO: \CursoCaratula\par}
\vspace{.30in}{\fontsize{14}{17}\selectfont\bfseries\color{modelblue!80!black}\TituloCaratula\par}
\vspace{.10in}{\fontsize{11}{13}\selectfont\bfseries\UbicacionCaratula\par}
\vspace{.35in}{\fontsize{13}{16}\selectfont\bfseries\TipoDocumentoCaratula\par}
\vspace{.30in}\begin{flushleft}\fontsize{10}{13}\selectfont\renewcommand{\arraystretch}{1.2}
\begin{tabular}{>{\bfseries}p{1.8in}p{3.05in}}\FichaCaratula\end{tabular}\end{flushleft}
\vfill{\fontsize{10}{12}\selectfont PUNO -- PERÚ\par}{\fontsize{9}{11}\selectfont 2026\par}
\end{minipage}\end{titlepage}

\clearpage\hypersetup{pageanchor=true}\pagenumbering{arabic}\tableofcontents\clearpage
\chapter*{ESPECIFICACIONES Y MONTAJE}\addcontentsline{toc}{chapter}{ESPECIFICACIONES Y MONTAJE}
metrado\_\allowbreak{}detallado.csv conserva ID, recurso, fórmula, poste/vano y estado. Los civiles no reciben volumen cerrado: excavación=\ensuremath{\pi}D\textsuperscript{2}h/4 y concreto=excavación-volumen desplazado cuando se defina geometría real. Un bloque 1.2 m o una malla dibujada no acredita capacidad ni resistividad. Las planillas dejan esas dimensiones vacías.


\setcounter{section}{6}
\section{Materiales, montaje y pruebas}
Postes: solicitar carga transversal de trabajo/rotura, posición del ensayo, dimensiones, peso, capacidad axial, interacciónN-M, comportamiento de fisuración, certificados por lote, orificios y transporte. Conductores: ficha AAAC, R20, masa, diámetro, GMR, RTS, E, \ensuremath{\alpha}, curva tensión/deformación/fluencia y carga mínima de grapas. Aisladores/crucetas: comprobar cargas con signo, distancias y fijación CAC de cada función.

Transformadores/celdas: capacidad neta a altitud local, relaciones/tomas, uk, pérdidas, grupo, neutro, aislamiento y ensayos. Montaje: replantear y liberar predios; ejecutar geotecnia/cimentaciones; izar con procedimiento del proveedor; tender por poleas; ajustar tabla de tensado a temperatura medida; comprobar puesta a tierra/protecciones y energizar tras pruebas. La tabla de tensado de esta revisión es orientativa para el modelo, no orden de tendido.


\setcounter{section}{7}
\section{Presupuesto, cronograma y cierre}
Se reemplaza el total aparente cerrado por una plantilla con cantidades trazables y precios sin cotizar. El presupuestoR01 se conserva como antecedente estimado; no se reutiliza como oferta local. Suministro, montaje, transporte, obras civiles y gestión se separan para evitar doble contabilización. Se solicita precio entregado en destino y tratamiento de IGV por oferta. APU, gastos, utilidad y cronograma de rendimiento permanecen editables.

Secuencia propuesta:1) confirmar curso y fuente;2) verificar generación y demanda;3) comparar tensión con cotizaciones;4) levantar corredor/cruces y predios;5) resolver mecánica/aislamiento/protección/tierra/cimentación;6) congelar planillas;7) generar CAD, metrados y presupuesto;8) revisión de responsable y pruebas. Duraciones en hoja de planificación son hipótesis, sin fecha de obra comprometida.


\end{document}
